% !PRINTED 167
\weberpage{167}

Wir haben im vorigen Abschnitt gesehen, daß, wenn

\weq{3}{\varkappa^2 = \dfrac{\vartheta_{10}^4}{\vartheta_{00}^4}}

ist, diese Aufgabe durch die elliptischen Funktionen gelöst wird, und zwar in der Weise, daß $x$, $y$, $z$ in der ganzen $v$-Ebene eindeutige und, außer wo sie unendlich sind, stetige Funktionen von $v$ werden. Ist nun aber $\varkappa^2$ gegeben, so entsteht die Frage, ob es zu jedem Wert von $\varkappa^2$ einen der Bedingung (3) genügenden Wert von $\omega$ gibt und ob es mehrere solche gibt.

Wir beweisen zunächst, daß durch die Differentialgleichungen (1) mit den Nebenbedingungen die Funktionen $x$, $y$, $z$ \gesperrt{eindeutig} bestimmt sind.

Es ergibt sich zunächst aus (1), daß $x^2+y^2$, $\varkappa^2x^2+z^2$ konstant sind, und aus den Nebenbedingungen folgt:

\weq{4}{y^2 = 1-x^2, \qquad z^2 = 1-\varkappa^2x^2.}

Angenommen, es existiere ein zweites System denselben Bedingungen genügender Funktionen $x_1$, $y_1$, $z_1$, so ist zunächst

\[
y_1^2 = 1-x_1^2, \qquad z_1^2 = 1-\varkappa^2x_1^2,
\]

und eine einfache Differentiation ergibt, daß

\[
\frac{x\,y_1z_1 - x_1\,yz}{1-\varkappa^2x^2x_1^2}
\]

konstant ist. (Das Additionstheorem der elliptischen Integrale, aus dem man die Form dieses Ausdruckes erhält, wird unmittelbar durch Differentiation bestätigt.) Aus den Nebenbedingungen folgt aber, daß diese Konstante verschwindet, also:

\[
x\,y_1z_1 = x_1\,yz,
\]

woraus man leicht schließt, indem man beiderseits quadriert:

\[
x = x_1, \qquad y=y_1, \qquad z=z_1.
\]

\begin{center}
\textbf{\S\,52. \gesperrt{Die unabhängige Variable $\varkappa^2$. Lineare Differentialgleichung für $K$.}}
\end{center}

Wir zeigen nun zunächst, daß und wie man zu einem beliebig gegebenen Wert von $\varkappa^2$ wenigstens \gesperrt{einen} der Bedingung

\weq{1}{\varkappa^2 = \dfrac{\vartheta_{10}^4}{\vartheta_{00}^4}}

genügenden Wert von $\omega$ finden kann. Diese Frage wird am vollständigsten beantwortet durch Zurückführung auf eine lineare Differentialgleichung zweiter Ordnung.
